\documentclass[10pt,letterpaper]{article}

\usepackage[margin=0.5in]{geometry}
\usepackage[T1]{fontenc}
\usepackage[utf8]{inputenc}
\usepackage{lmodern}
\usepackage[hidelinks]{hyperref}
\usepackage{enumitem}
\usepackage{titlesec}

\pagestyle{empty}
\setlength{\parindent}{0pt}
\setlength{\parskip}{0pt}
\setlist[itemize]{leftmargin=12pt, itemsep=1pt, topsep=1pt, parsep=0pt, partopsep=0pt}

\titleformat{\section}{\large\bfseries}{}{0pt}{}[\titlerule]
\titlespacing*{\section}{0pt}{5pt}{2pt}

\newcommand{\resumeSubheading}[4]{%
  \textbf{#1} \hfill #2 \\
  #3 \hfill #4 \\
}

\begin{document}

\begin{center}
    {\Large \textbf{Sanskar Jain}} \\
    Full-Stack Engineer \;|\; C\# / .NET \;|\; Agentic AI \;|\; LLM Integrations \\
    Hyderabad, India \enspace | \enspace 9021620156 \enspace | \enspace
    \href{mailto:sanskarjain0901@gmail.com}{sanskarjain0901@gmail.com} \enspace | \enspace
    \href{https://www.linkedin.com/in/sanskar-jain0901/}{LinkedIn}
\end{center}

\vspace{2pt}

\section*{Summary}
Full-stack software engineer with 2+ years building \textbf{C\#/.NET} (ASP.NET Core) REST APIs and microservices, plus hands-on experience developing \textbf{LLM-powered internal agents} using the \textbf{OpenAI API}, prompt engineering, and custom agent orchestration at Verisk Analytics. Comfortable across the full stack --- relational/\textbf{NoSQL} databases, JavaScript/HTML/CSS front-ends, and unit-tested, Agile-delivered backend services --- with strong interest in applying Generative AI to real business workflows.

\section*{Skills}
\textbf{Agentic AI \& GenAI:} OpenAI API, Prompt Engineering, Context Management, Agent/Workflow Orchestration, LLM Integration Patterns \\
\textbf{Languages \& Web:} C\#, JavaScript, TypeScript, HTML5, CSS3, JSON, Python \\
\textbf{Frameworks \& Platforms:} .NET 8, ASP.NET Core, ASP.NET MVC, Web API, REST APIs, Microservices, Kafka, Redis \\
\textbf{Databases:} SQL Server, MySQL, PostgreSQL, MongoDB (NoSQL), Query Optimization, Database Design \\
\textbf{Testing \& Tools:} xUnit, NUnit, Git, Azure DevOps, Swagger/OpenAPI, Postman, Docker \\
\textbf{Practices:} Agile/Scrum, RESTful API Design, Code Reviews, CI/CD, Technical Documentation

\section*{Experience}
\resumeSubheading{Verisk Analytics}{Jan 2024 -- Present}{Software Engineer II}{Hyderabad, India}
\begin{itemize}
    \item Built \textbf{LLM-powered internal workflow agents} using the \textbf{OpenAI API}, integrating prompt-driven automation for \textbf{Jira ticket creation} and task routing, reducing manual coordination effort and speeding up execution across teams.
    \item Designed \textbf{prompt templates and context-passing logic} to reliably map business rules and validation checks into agent-driven actions, connecting LLM outputs with backend services via structured \textbf{REST/JSON} contracts.
    \item Built \textbf{ASP.NET Core Web API} endpoints handling \textbf{200K+ weekly requests}, with structured error handling, input validation, and \textbf{Swagger/OpenAPI} documentation across a full-stack \textbf{C\#/.NET} application.
    \item Co-designed \textbf{RulesRunner}, a stateless \textbf{C\#} microservice using \textbf{CQRS} and SQS, reducing client wait time from \textbf{3000ms to 500ms (83\%)} and enabling horizontal scaling for compliance workloads.
    \item Engineered a scalable import/export backend for \textbf{Xactimate} across \textbf{30+ data types} using async \textbf{C\#} pipelines, \textbf{MongoDB} (NoSQL), and \textbf{SQL}-backed storage, cutting transfer time by \textbf{38\%} and issues by \textbf{25\%}.
    \item Improved \textbf{API response time by 34\%} via \textbf{Redis} caching and \textbf{SQL query optimization}; wrote \textbf{xUnit/NUnit} unit and integration tests and participated in structured \textbf{code reviews} across the team.
    \item Overhauled \textbf{Azure DevOps CI/CD} pipelines with staged gates and parallel jobs, cutting deployment time by \textbf{50\%} and build cycles by \textbf{33\%} within an \textbf{Agile} delivery model.
    \item Reduced technical debt by \textbf{55\%} and improved runtime performance by \textbf{30\%} by migrating legacy services to \textbf{.NET 8}, staying current with modern async patterns and new tooling.
\end{itemize}

\section*{Projects}

\resumeSubheading{Internal Workflow Automation Agent}{2025 -- 2026}{}{OpenAI API, C\#/Python, Prompt Engineering, Jira, REST APIs}
\begin{itemize}
    \item Built an \textbf{LLM-based agent} that parses task requests and drafts/creates \textbf{Jira tickets} via prompt-engineered calls to the \textbf{OpenAI API}, with custom orchestration logic for context management and input validation.
    \item Connected agent output to downstream backend services and business rules, reducing repetitive manual coordination and improving turnaround time for operational requests.
\end{itemize}

\resumeSubheading{Expense Tracker -- Full-Stack Microservices Application}{Aug 2023 -- Dec 2023}{}{C\#, Spring Boot, JavaScript, HTML/CSS, Kafka, SQL/NoSQL}
\begin{itemize}
    \item Built \textbf{5+ microservices} with async \textbf{Kafka} messaging, a \textbf{JavaScript}/HTML/CSS UI layer consuming \textbf{REST/JSON} APIs, and \textbf{relational and NoSQL} data persistence, containerised with \textbf{Docker}.
\end{itemize}

\section*{Certifications}
\textbf{AWS Certified Cloud Practitioner} \hfill Amazon Web Services \;|\; Oct 2024

\section*{Education}
\textbf{B.Tech in Computer Engineering} --- Vishwakarma Institute of Technology, Pune (SPPU) \hfill 2024 \;|\; CGPA: 8.8

\end{document}
